\chapter{RSA}

\section{Die andere Aufgabe}

\begin{merksatz}
\boxtitel{Was Kapitel 12 nicht leistet}
Der Schlüsseltausch verlangt, dass beide Seiten gleichzeitig anwesend sind und miteinander sprechen.

\begin{center}
{\small
\renewcommand{\arraystretch}{1.5}
\begin{tabular}{@{}P{2.4cm}P{6.2cm}@{}}
\toprule
\textbf{Nicht möglich} & Einem Empfänger eine Nachricht senden, der gerade nicht erreichbar ist.\\
\textbf{Nicht möglich} & Eine Nachricht so unterschreiben, dass jeder die Echtheit prüfen kann.\\
\bottomrule
\end{tabular}}
\end{center}

\begin{center}
\begin{minipage}{0.9\textwidth}
\textbf{Gesucht ist ein Schlüsselpaar: einer zum Verschließen, ein anderer zum Öffnen.} Der erste darf öffentlich sein.
\end{minipage}
\end{center}

\schritt{Der Leitfall}
Eine Praxis sendet einen Arztbrief an das Klinikum. Sie hat noch nie mit ihm gesprochen, kennt aber dessen öffentlichen Schlüssel.

\begin{center}
\begin{minipage}{0.88\textwidth}
\kursiv{Diffie und Hellman hatten die Idee 1976 formuliert, ohne ein Verfahren angeben zu können. \textbf{Ein Jahr später fanden drei Mathematiker am MIT eines} --- und es beruht auf einer Frage, die zweitausend Jahre alt ist: wie schwer es ist, eine Zahl in ihre Primfaktoren zu zerlegen.}
\end{minipage}
\end{center}
\end{merksatz}

\section{Der Satz von Euler}

\begin{merksatz}
\boxtitel{Die Eulersche Funktion}
Für eine natürliche Zahl $n$ ist $\varphi(n)$ die Anzahl der zu $n$ teilerfremden Zahlen zwischen $1$ und $n$.

\begin{center}
{\small
\renewcommand{\arraystretch}{1.5}
\begin{tabular}{@{}P{2.4cm}P{2.6cm}P{3.0cm}@{}}
\toprule
$n$ & $\varphi(n)$ & \\
\midrule
$p$ prim & $p-1$ & alle außer $0$\\
$p\cdot q$ & $(p-1)(q-1)$ & für $p \ne q$ prim\\
$143 = 11\cdot13$ & $10\cdot12 = 120$ & \\
\bottomrule
\end{tabular}}
\end{center}

\schritt{Warum $(p-1)(q-1)$}
Von den $n = pq$ Zahlen sind genau die Vielfachen von $p$ und die von $q$ nicht teilerfremd:
\[ n - q - p + 1 = pq-p-q+1 = (p-1)(q-1) \]
Die Eins wird addiert, weil die Null in beiden Listen steht. \checkmark
\end{merksatz}

\begin{merksatz}
\boxtitel{Satz 13.1 --- der Satz von Euler}
Ist $\gcd(a,n) = 1$, so gilt
\[ a^{\varphi(n)} \equiv 1 \pmod n . \]

\schritt{Der Beweis in Kürze}
Die zu $n$ teilerfremden Reste bilden eine Gruppe mit $\varphi(n)$ Elementen. Multipliziert man alle mit $a$, so entsteht dieselbe Menge in anderer Reihenfolge --- denn $a$ ist umkehrbar.

Das Produkt aller Elemente ist damit einerseits $P$, andererseits $a^{\varphi(n)}P$. Kürzen liefert die Behauptung. $\blacksquare$

\begin{center}
\begin{minipage}{0.88\textwidth}
\kursiv{Der kleine Satz von Fermat aus Kapitel~12 ist der Sonderfall $n=p$. Der vollständige Beweis steht in Band 6, Teil~II --- er ist der Satz von Lagrange über Gruppenordnungen, auf diese Gruppe angewandt.}
\end{minipage}
\end{center}
\end{merksatz}

\section{Das Verfahren}

\begin{merksatz}
\boxtitel{Die Schlüsselerzeugung}
\begin{center}
{\small
\renewcommand{\arraystretch}{1.5}
\begin{tabular}{@{}P{1.0cm}P{7.4cm}@{}}
\toprule
$1$ & Wähle zwei große Primzahlen $p$ und $q$ und setze $n = pq$.\\
$2$ & Berechne $\varphi(n) = (p-1)(q-1)$.\\
$3$ & Wähle $e$ teilerfremd zu $\varphi(n)$.\\
$4$ & Bestimme $d$ mit $e\,d \equiv 1 \pmod{\varphi(n)}$.\\
\bottomrule
\end{tabular}}
\end{center}

\begin{center}
\renewcommand{\arraystretch}{1.5}
\begin{tabular}{@{}P{2.6cm}P{2.4cm}P{2.8cm}@{}}
\toprule
& \textbf{Inhalt} & \textbf{Status}\\
\midrule
Öffentlicher Schlüssel & $(n,e)$ & \textcolor{gruen}{veröffentlicht}\\
Geheimer Schlüssel & $d$ & \textcolor{akzent}{geheim}\\
Zu vernichten & $p$, $q$, $\varphi(n)$ & \textcolor{akzent}{geheim}\\
\bottomrule
\end{tabular}
\end{center}

\schritt{Die Operationen}
\[ c = m^{e} \bmod n, \qquad m = c^{d} \bmod n \]
\end{merksatz}

\begin{merksatz}
\boxtitel{Satz 13.2 --- warum es funktioniert}
Für alle $m$ mit $\gcd(m,n)=1$ gilt
\[ \bigl(m^{e}\bigr)^{d} \equiv m \pmod n . \]

\schritt{Beweis}
Aus $e\,d \equiv 1 \pmod{\varphi(n)}$ folgt
\[ e\,d = 1 + k\,\varphi(n) \quad\text{für ein } k\in\mathbb{N} . \]

Damit:
\begin{align*}
m^{ed} &= m^{\,1+k\varphi(n)}\\
&= m\cdot\bigl(m^{\varphi(n)}\bigr)^{k}\\
&\equiv m\cdot 1^{k} && \text{(Satz 13.1)}\\
&= m
\end{align*}
$\blacksquare$

\begin{center}
\begin{minipage}{0.88\textwidth}
\kursiv{Der Beweis ist vier Zeilen lang und benutzt nur den Satz von Euler. \textbf{Die Aussage gilt tatsächlich für alle $m$}, auch für die nicht teilerfremden --- das erfordert eine zusätzliche Betrachtung über den chinesischen Restsatz. Bei großen Primzahlen ist der Fall ohnehin verschwindend selten: Ein $m$ mit $\gcd(m,n)>1$ zu finden hieße, $n$ zu faktorisieren.}
\end{minipage}
\end{center}
\end{merksatz}

\begin{merksatz}
\boxtitel{Ein vollständiges Beispiel}
\begin{center}
{\small
\renewcommand{\arraystretch}{1.48}
\begin{tabular}{@{}P{2.2cm}P{2.4cm}P{2.8cm}@{}}
\toprule
\textbf{Größe} & \textbf{Wert} & \textbf{Rechnung}\\
\midrule
$p$, $q$ & $11$, $13$ & gewählt\\
$n$ & $143$ & $11\cdot13$\\
$\varphi(n)$ & $120$ & $10\cdot12$\\
$e$ & $7$ & teilerfremd zu $120$\\
$d$ & $\mathbf{103}$ & $7\cdot103 = 721$\\
\bottomrule
\end{tabular}}
\end{center}

\kursiv{Probe für $d$:} $721 = 6\cdot120+1$ \checkmark

\schritt{Verschlüsseln von $m = 9$}
\[ c = 9^{7} \bmod 143 \]

\kursiv{Quadrieren und Multiplizieren}, $7 = \texttt{111}_2$:
\begin{center}
\renewcommand{\arraystretch}{1.46}
\begin{tabular}{@{}P{1.2cm}P{2.6cm}P{2.4cm}@{}}
\toprule
\textbf{Bit} & \textbf{Rechnung} & \textbf{Ergebnis}\\
\midrule
$1$ & $1^2\cdot9$ & $9$\\
$1$ & $9^2\cdot9 = 729$ & $14$\\
$1$ & $14^2\cdot9 = 1764$ & $\mathbf{48}$\\
\bottomrule
\end{tabular}
\end{center}

\kursiv{Proben:} $729 = 5\cdot143+14$ \checkmark; $1764 = 12\cdot143+48$ \checkmark

\schritt{Entschlüsseln}
\[ 48^{103} \bmod 143 = 9 \quad \checkmark \]

\begin{center}
\begin{minipage}{0.88\textwidth}
\kursiv{Diese Rechnung ist von Hand mühsam --- zehn Quadrierungen. Der Empfänger kennt allerdings $p$ und $q$ und kann über den chinesischen Restsatz modulo $11$ und modulo $13$ getrennt rechnen; das ist etwa viermal schneller. \textbf{Der geheime Schlüssel enthält mehr als nur $d$}, und praktische Umsetzungen speichern $p$, $q$ und Hilfsgrößen mit.}
\end{minipage}
\end{center}
\end{merksatz}

\section{Wie sicher ist es}

\begin{merksatz}
\boxtitel{Wovon die Sicherheit abhängt}
\begin{center}
{\small
\renewcommand{\arraystretch}{1.52}
\begin{tabular}{@{}P{2.8cm}P{5.8cm}@{}}
\toprule
\textbf{Faktorisierung} & Kennt man $p$ und $q$, so kennt man $\varphi(n)$ und damit $d$.\\
\textbf{$\varphi(n)$ bestimmen} & Gleichwertig zur Faktorisierung: Aus $n$ und $\varphi(n)$ folgen $p$ und $q$ über eine quadratische Gleichung.\\
\textbf{$d$ bestimmen} & Ebenfalls gleichwertig --- daraus lässt sich $n$ faktorisieren.\\
\textbf{RSA-Problem} & Aus $c$ die $e$-te Wurzel ziehen. \textcolor{mittel}{Höchstens} so schwer wie Faktorisieren.\\
\bottomrule
\end{tabular}}
\end{center}

\begin{center}
\begin{minipage}{0.88\textwidth}
\kursiv{Die letzte Zeile hat denselben Vorbehalt wie beim Diffie-Hellman-Problem in Kapitel~12: Ob das Wurzelziehen \kursiv{genauso} schwer ist wie das Faktorisieren, ist offen. \textbf{Man weiß nur, dass es nicht schwerer ist.}}
\end{minipage}
\end{center}

\schritt{Der Stand der Faktorisierung}
\begin{center}
{\small
\renewcommand{\arraystretch}{1.48}
\begin{tabular}{@{}P{2.2cm}P{1.8cm}P{2.0cm}P{2.0cm}@{}}
\toprule
\textbf{Zahl} & \textbf{Bits} & \textbf{Jahr} & \textbf{Aufwand}\\
\midrule
RSA-768 & $768$ & $2009$ & $\approx2000$ Kernjahre\\
RSA-250 & $829$ & $2020$ & $\approx2700$ Kernjahre\\
$2048$ Bit & $2048$ & --- & \textcolor{gruen}{unerreicht}\\
\bottomrule
\end{tabular}}
\end{center}

\begin{center}
{\small
\renewcommand{\arraystretch}{1.48}
\begin{tabular}{@{}P{2.2cm}P{2.2cm}P{2.8cm}@{}}
\toprule
$n$ & \textbf{Aufwand} & \\
\midrule
$2^{768}$ & $2^{76{,}5}$ & gebrochen\\
$2^{1024}$ & $2^{86{,}8}$ & \textcolor{mittel}{grenzwertig}\\
$2^{2048}$ & $2^{116{,}9}$ & sicher\\
$2^{3072}$ & $2^{138{,}7}$ & sicher\\
\bottomrule
\end{tabular}}
\end{center}

\begin{center}
\begin{minipage}{0.88\textwidth}
\kursiv{Es sind dieselben Zahlen wie beim diskreten Logarithmus --- das Zahlkörpersieb löst beide Probleme mit demselben Aufwandsgesetz. \textbf{Die Schlüssellängenempfehlungen für RSA und Diffie-Hellman stimmen deshalb überein}, und ein Fortschritt bei einem Problem betrifft meist auch das andere.}
\end{minipage}
\end{center}
\end{merksatz}

\section{Der Bruchfall}

\begin{warnung}
\boxtitelr{Das reine Verfahren ist unsicher}
Die Vorschrift $c = m^{e}$ --- ohne weitere Maßnahmen --- hat mehrere schwere Mängel.

\begin{center}
{\small
\renewcommand{\arraystretch}{1.52}
\begin{tabular}{@{}P{2.4cm}P{6.2cm}@{}}
\toprule
\textbf{Bestimmt} & Dieselbe Nachricht ergibt stets denselben Geheimtext. Bei wenigen möglichen Nachrichten rät man sie durch Ausprobieren.\\
\textbf{Formbar} & $c_1\cdot c_2$ entschlüsselt zu $m_1\cdot m_2$.\\
\textbf{Kleine Werte} & Ist $m^{e} < n$, so entfällt die Modulrechnung --- man zieht einfach die $e$-te Wurzel.\\
\bottomrule
\end{tabular}}
\end{center}

\schritt{Zur ersten Zeile}
Ein Arztbrief mit dem Feld \enquote{Befund: positiv} oder \enquote{negativ} hat zwei mögliche Werte. Der Angreifer verschlüsselt beide selbst --- er kennt ja den öffentlichen Schlüssel --- und vergleicht.

\begin{center}
\begin{minipage}{0.9\textwidth}
\textbf{Ein Verfahren mit öffentlichem Schlüssel, das deterministisch ist, kann nicht sicher sein.} Der Angreifer kann jede Vermutung selbst prüfen.
\end{minipage}
\end{center}

\schritt{Zur zweiten Zeile, nachgerechnet}
Mit $n=143$, $e=7$:
\begin{center}
\renewcommand{\arraystretch}{1.46}
\begin{tabular}{@{}P{2.0cm}P{2.0cm}P{2.4cm}@{}}
\toprule
$m$ & $c$ & \\
\midrule
$5$ & $47$ & \\
$7$ & $6$ & \\
\midrule
$35$ & $\mathbf{139}$ & $47\cdot6 = 282 \equiv 139$\\
\bottomrule
\end{tabular}
\end{center}

\kursiv{Probe:} $282 = 143+139$ \checkmark, und $35^{7} \bmod 143 = 139$ \checkmark

\begin{center}
\begin{minipage}{0.88\textwidth}
\kursiv{Der Angreifer kann also aus dem Geheimtext einer Zahl den Geheimtext ihres Vielfachen erzeugen --- ohne den Klartext zu kennen. Das ist dieselbe Formbarkeit wie beim Einmalschlüssel in Kapitel~9, nur multiplikativ statt additiv.}
\end{minipage}
\end{center}

\schritt{Die Abhilfe}
Man füllt die Nachricht vor dem Potenzieren mit Zufall auf, nach einem genau festgelegten Verfahren.

\begin{center}
\begin{minipage}{0.88\textwidth}
\kursiv{Damit ist die Verschlüsselung nicht mehr bestimmt --- dieselbe Nachricht ergibt jedes Mal einen anderen Geheimtext --- und die Struktur des Auffüllens zerstört die Formbarkeit. \textbf{Die Auffüllung ist kein technisches Beiwerk, sondern der Teil, der das Verfahren sicher macht.} Frühere Auffüllverfahren waren angreifbar, und der bekannteste dieser Angriffe nutzte aus, dass der Empfänger fehlerhafte Auffüllungen meldete.}
\end{minipage}
\end{center}
\end{warnung}

\begin{warnung}
\boxtitelr{Und vier Umsetzungsfallen}
\begin{center}
{\small
\renewcommand{\arraystretch}{1.5}
\begin{tabular}{@{}P{2.4cm}P{6.2cm}@{}}
\toprule
\textbf{Gemeinsame Faktoren} & Werden Primzahlen mit schlechtem Zufall erzeugt, so teilen sich verschiedene Schlüssel einen Faktor --- und beide fallen.\\
\textbf{Zu nahe Primzahlen} & Liegen $p$ und $q$ dicht beieinander, so findet man sie durch Suche um $\sqrt n$.\\
\textbf{Laufzeit} & Die Entschlüsselungsdauer hängt von $d$ ab und verrät es.\\
\textbf{Fehlerangriff} & Ein einziger Rechenfehler bei der Rechnung über den chinesischen Restsatz verrät $p$.\\
\bottomrule
\end{tabular}}
\end{center}

\schritt{Zur ersten Zeile}
Eine Untersuchung von 2012 prüfte Millionen öffentlicher Schlüssel aus dem Internet und berechnete paarweise größte gemeinsame Teiler.

\begin{center}
\begin{minipage}{0.88\textwidth}
\kursiv{Bei einigen Zehntausend Schlüsseln wurde man fündig --- meist von eingebetteten Geräten, die ihre Schlüssel kurz nach dem Einschalten erzeugten, bevor genug Zufall gesammelt war. \textbf{Der größte gemeinsame Teiler ist in Millisekunden berechnet}, während die Faktorisierung einer einzelnen Zahl aussichtslos ist. Zwei schlechte Schlüssel brechen einander.}
\end{minipage}
\end{center}

\schritt{Zur vierten Zeile}
Rechnet der Empfänger getrennt modulo $p$ und modulo $q$ und tritt bei einer der beiden Teilrechnungen ein Fehler auf, so gilt für die fehlerhafte Signatur $s'$:
\[ \gcd\bigl(s'^{\,e} - m,\ n\bigr) = q \]

\begin{center}
\begin{minipage}{0.88\textwidth}
\kursiv{Ein einziger gekippter Bitwert --- durch Spannungsschwankung, Hitze oder gezielte Einstrahlung --- gibt den Schlüssel preis. Deshalb prüfen sichere Umsetzungen jede Signatur nach dem Erzeugen. \textbf{Ein mathematisch einwandfreies Verfahren kann durch einen Hardwarefehler vollständig fallen} --- und das ist ein Angriffsweg, den kein Sicherheitsbeweis erfasst.}
\end{minipage}
\end{center}
\end{warnung}

\section{Übungsaufgaben}

\begin{aufgabe}
\boxtitel{Aufgabe 13.1 --- RSA durchrechnen}
Mit $p = 5$, $q = 11$ und $e = 3$: Bestimmen Sie $d$ und verschlüsseln Sie $m = 8$.

\loesung
\schritt{Schritt 1 --- Die Kenngrößen}
\[ n = 55, \qquad \varphi(n) = 4\cdot10 = 40 \]

\kursiv{Teilerfremdheit prüfen:} $\gcd(3,40) = 1$ \checkmark

\schritt{Schritt 2 --- $d$ bestimmen}
Gesucht ist $d$ mit $3d \equiv 1 \pmod{40}$.

\kursiv{Durchprobieren der Vielfachen:} $3\cdot27 = 81 = 2\cdot40+1$ \checkmark
\[ d = 27 \]

\schritt{Schritt 3 --- Verschlüsseln}
\[ c = 8^{3} = 512 \]
\[ 512 = 9\cdot55 + 17 \quad\Longrightarrow\quad c = \mathbf{17} \]

\schritt{Schritt 4 --- Entschlüsseln}
\[ 17^{27} \bmod 55 \]

\kursiv{Schrittweise:} $17^2 = 289 = 5\cdot55+14$, also $14$.
\[ 17^{4} = 14^{2} = 196 = 3\cdot55+31 \ \to\ 31 \]
\[ 17^{8} = 31^{2} = 961 = 17\cdot55+26 \ \to\ 26 \]
\[ 17^{16} = 26^{2} = 676 = 12\cdot55+16 \ \to\ 16 \]

Mit $27 = 16+8+2+1$:
\[ 17^{27} = 16\cdot26\cdot14\cdot17 \]
\[ 16\cdot26 = 416 = 7\cdot55+31 \ \to\ 31 \]
\[ 31\cdot14 = 434 = 7\cdot55+49 \ \to\ 49 \]
\[ 49\cdot17 = 833 = 15\cdot55+8 \ \to\ \mathbf{8} \quad \checkmark \]

\schritt{Schritt 5 --- Die Probe über Satz 13.2}
\[ e\,d = 81 = 2\cdot40+1 \]
\[ m^{81} = m\cdot(m^{40})^{2} \equiv m\cdot1 = m \]

\begin{center}
\begin{minipage}{0.86\textwidth}
\kursiv{Zur Anwendbarkeit von Satz 13.1: $\gcd(8,55) = 1$ \checkmark. Bei $m = 5$ oder $m=11$ wäre der Satz nicht unmittelbar anwendbar --- die Aussage gilt dort trotzdem, wie die erweiterte Betrachtung zeigt.}
\end{minipage}
\end{center}

\schritt{Schritt 6 --- Warum dieses Beispiel unsicher ist}
Mit $n = 55$ gibt es $55$ mögliche Nachrichten. Der Angreifer verschlüsselt alle und vergleicht.

\begin{center}
\begin{minipage}{0.88\textwidth}
\kursiv{Und er kann $55 = 5\cdot11$ im Kopf faktorisieren. \textbf{Das Beispiel zeigt die Rechnung; die Sicherheit kommt allein aus der Größe} --- praktisch verwendet man Primzahlen mit über tausend Bit.}
\end{minipage}
\end{center}
\end{aufgabe}

\begin{aufgabe}
\boxtitel{Aufgabe 13.2 --- Die Formbarkeit ausnutzen}
Ein Angreifer kennt den öffentlichen Schlüssel $(n,e)$ und einen Geheimtext $c$, dessen Klartext $m$ er erfahren will. Er darf den Empfänger \kursiv{einen} anderen Geheimtext entschlüsseln lassen --- aber nicht $c$ selbst.

\loesung
\schritt{Schritt 1 --- Der Ansatz}
Der Angreifer wählt ein $r$ und bildet
\[ c' = c\cdot r^{e} \bmod n . \]

\schritt{Schritt 2 --- Was der Empfänger zurückgibt}
\[ (c')^{d} = (c\cdot r^{e})^{d} = c^{d}\cdot r^{ed} \equiv m\cdot r \pmod n \]

\schritt{Schritt 3 --- Das Ergebnis}
Der Angreifer erhält $m\cdot r$ und dividiert:
\[ m = (m\cdot r)\cdot r^{-1} \bmod n \]

\begin{center}
\begin{minipage}{0.9\textwidth}
\textbf{Er kennt den Klartext --- und $c'$ sah für den Empfänger völlig unverdächtig aus.}
\end{minipage}
\end{center}

\schritt{Schritt 4 --- Ein Zahlenbeispiel}
$n=55$, $e=3$, $c=17$ (also $m=8$). Der Angreifer wählt $r=2$:
\[ c' = 17\cdot2^{3} = 17\cdot8 = 136 = 2\cdot55+26 \ \to\ 26 \]

Der Empfänger entschlüsselt: $26^{27} \bmod 55 = 16$.
\[ m = 16\cdot2^{-1} \bmod 55 \]

\kursiv{Inverses von $2$:} $2\cdot28 = 56 \equiv 1$, also $2^{-1} = 28$.
\[ 16\cdot28 = 448 = 8\cdot55+8 \ \to\ \mathbf{8} \quad \checkmark \]

\schritt{Schritt 5 --- Die Deutung}
\begin{center}
\begin{minipage}{0.88\textwidth}
\kursiv{Dieser Angriff setzt voraus, dass der Empfänger überhaupt bereit ist, fremde Geheimtexte zu entschlüsseln. Das klingt weltfremd und ist es nicht: Ein Server, der auf fehlerhafte Nachrichten mit einer Fehlermeldung antwortet, gibt bereits genug preis. \textbf{Jede Antwort auf einen gewählten Geheimtext ist eine Entschlüsselung} --- und sei sie nur ein Ja oder Nein.}
\end{minipage}
\end{center}

\schritt{Schritt 6 --- Die Abhilfe}
Eine ordentliche Auffüllung macht $c'$ mit überwältigender Wahrscheinlichkeit ungültig.

\begin{center}
\begin{minipage}{0.88\textwidth}
\kursiv{Das Produkt $m\cdot r$ hat keine gültige Auffüllstruktur, und der Empfänger lehnt ab. Damit ist der Angriff abgewehrt --- vorausgesetzt, die Ablehnung verrät nicht, \kursiv{warum} abgelehnt wurde, und die Antwortzeit ist in allen Fällen gleich. \textbf{Die Abwehr liegt außerhalb der Mathematik}, und genau dort versagte sie historisch mehrfach.}
\end{minipage}
\end{center}
\end{aufgabe}

\begin{aufgabe}
\boxtitel{Aufgabe 13.3 --- Signaturen}
Erklären Sie, wie sich RSA zum Unterschreiben verwenden lässt, und warum man nicht die Nachricht selbst signiert.

\loesung
\schritt{Schritt 1 --- Die Vertauschung}
Beim Verschlüsseln wendet man $e$ an und der Empfänger $d$. Beim Signieren ist es umgekehrt:
\[ s = m^{d} \bmod n \qquad\text{(nur der Inhaber kann das)} \]
\[ s^{e} \stackrel{?}{=} m \qquad\text{(jeder kann das prüfen)} \]

\begin{center}
\begin{minipage}{0.9\textwidth}
\textbf{Das funktioniert, weil $e$ und $d$ symmetrisch zueinander stehen:} $m^{ed} = m^{de}$.
\end{minipage}
\end{center}

\schritt{Schritt 2 --- Was das leistet}
\begin{center}
{\small
\renewcommand{\arraystretch}{1.5}
\begin{tabular}{@{}P{3.2cm}P{5.4cm}@{}}
\toprule
\textbf{Herkunft} & Nur wer $d$ kennt, kann $s$ erzeugen.\\
\textbf{Unversehrtheit} & Jede Änderung an $m$ macht die Prüfung ungültig.\\
\textbf{Nicht\-abstreitbarkeit} & Anders als beim Beglaubigungscode aus Kapitel~11 kann der Empfänger die Signatur nicht selbst erzeugt haben.\\
\bottomrule
\end{tabular}}
\end{center}

\schritt{Schritt 3 --- Warum man nicht $m$ selbst signiert}
\kursiv{Erstens: Länge.} Die Nachricht muss kleiner als $n$ sein. Ein Arztbrief ist es nicht.

\kursiv{Zweitens: Fälschbarkeit.} Der Angreifer wählt ein beliebiges $s$ und berechnet $m = s^{e}$.

\begin{center}
\begin{minipage}{0.9\textwidth}
\textbf{Er hat eine gültige Signatur --- für eine Nachricht, die er sich nicht aussuchen konnte.}
\end{minipage}
\end{center}

\kursiv{Drittens: Formbarkeit.} Aus zwei Signaturen $s_1$, $s_2$ ergibt sich
\[ s_1\cdot s_2 = (m_1 m_2)^{d} \]
--- eine gültige Signatur für $m_1 m_2$, ohne Kenntnis von $d$.

\schritt{Schritt 4 --- Die Abhilfe}
Man signiert den Hashwert:
\[ s = H(m)^{d} \bmod n \]

\begin{center}
{\small
\renewcommand{\arraystretch}{1.5}
\begin{tabular}{@{}P{2.4cm}P{6.2cm}@{}}
\toprule
\textbf{Länge} & Der Hashwert ist kurz genug.\\
\textbf{Fälschbarkeit} & Der Angreifer müsste zu $s^{e}$ ein passendes Urbild finden --- das ist Definition 11.1.\\
\textbf{Formbarkeit} & $H(m_1)\cdot H(m_2)$ ist kein Hashwert einer sinnvollen Nachricht.\\
\bottomrule
\end{tabular}}
\end{center}

\schritt{Schritt 5 --- Die Folge für die Hashfunktion}
\begin{center}
\begin{minipage}{0.9\textwidth}
\textbf{Findet jemand eine Kollision $H(m_1) = H(m_2)$, so überträgt sich jede Signatur von $m_1$ auf $m_2$.}
\end{minipage}
\end{center}

\begin{center}
\begin{minipage}{0.88\textwidth}
\kursiv{Deshalb war der Bruch von MD5 in Kapitel~11 folgenschwer: Man konnte zwei Zertifikate mit gleichem Hashwert und verschiedenem Inhalt bauen, eines harmlos signieren lassen und die Signatur auf das andere übertragen. \textbf{Die Sicherheit der Signatur ist durch die schwächste der beiden Komponenten begrenzt} --- und das war nicht RSA.}
\end{minipage}
\end{center}

\schritt{Schritt 6 --- Auch hier braucht es Auffüllung}
Die naheliegende Form $H(m)^{d}$ ohne weitere Struktur ist ebenfalls angreifbar.

\begin{center}
\begin{minipage}{0.86\textwidth}
\kursiv{Der Hashwert ist kurz, $n$ ist lang --- der Rest wird aufgefüllt, und die Art des Auffüllens war wiederholt Angriffspunkt. Ein bekannter Fall von 2006 nutzte aus, dass manche Prüfroutinen Zusatzbytes hinter dem Hashwert ignorierten; bei $e=3$ ließen sich damit Signaturen ohne Kenntnis von $d$ fälschen.}
\end{minipage}
\end{center}
\end{aufgabe}

\begin{geschichte}
\boxtitelg{Ein Wochenende im April}
\Person{Ron Rivest}, \Person{Adi Shamir} und \Person{Leonard Adleman} arbeiteten 1977 am MIT an der Frage, ob Diffie und Hellmans Idee eines Verfahrens mit öffentlichem Schlüssel überhaupt umsetzbar sei.

\begin{center}
\begin{minipage}{0.9\textwidth}
\kursiv{Die Arbeitsteilung war ungewöhnlich: Rivest und Shamir schlugen Verfahren vor, Adleman brach sie. Nach eigener Erinnerung etwa zweiundvierzigmal. Der dreiundvierzigste Vorschlag hielt stand.}
\end{minipage}
\end{center}

\schritt{Die Nacht der Entdeckung}
Rivest kam die Idee in der Nacht zum 3.\ April 1977, nach einem Abend mit Wein bei einem Pessachfest. Er schrieb die Arbeit noch in derselben Nacht.

\begin{center}
\begin{minipage}{0.9\textwidth}
\kursiv{Adleman wollte zunächst nicht als Autor genannt werden --- er habe nur Vorschläge widerlegt. Rivest bestand darauf. Die Reihenfolge der Namen ergab sich alphabetisch rückwärts, und so heißt das Verfahren RSA statt ARS.}
\end{minipage}
\end{center}

\schritt{Die Vorwegnahme}
\Person{Clifford Cocks} hatte dasselbe Verfahren 1973 beim britischen GCHQ gefunden --- in einer halben Stunde, wie er berichtete.

\begin{center}
\begin{minipage}{0.9\textwidth}
\kursiv{Er hielt es für rechnerisch unpraktikabel und schrieb es auf vier Seiten nieder, die bis 1997 geheim blieben. \textbf{Cocks hatte zudem keine Signaturen betrachtet} --- ihm ging es allein um Verschlüsselung. Der Gedanke, dass dasselbe Verfahren rückwärts eine Unterschrift liefert, fehlt in seiner Arbeit.}
\end{minipage}
\end{center}

\schritt{Die Zahlen der Preisaufgabe}
1977 veröffentlichte Rivest in einer Zeitschriftenkolumne eine mit $129$ Stellen verschlüsselte Nachricht und setzte hundert Dollar aus. Er schätzte den nötigen Aufwand auf $4\cdot10^{16}$ Jahre.

\begin{center}
\begin{minipage}{0.9\textwidth}
\kursiv{Die Zahl wurde 1994 faktorisiert --- von sechshundert Freiwilligen über das Internet, in acht Monaten. Der Klartext lautete \enquote{the magic words are squeamish ossifrage}. \textbf{Rivests Schätzung lag um vierzehn Größenordnungen daneben}, weil er die Entwicklung der Faktorisierungsverfahren nicht vorhersehen konnte. Das ist die deutlichste Warnung vor Sicherheitsprognosen, die dieses Fach kennt.}
\end{minipage}
\end{center}
\end{geschichte}

\begin{forschung}
\boxtitelp{RSA heute}
\begin{center}
{\small
\renewcommand{\arraystretch}{1.45}
\begin{tabular}{@{}P{2.8cm}P{5.4cm}@{}}
\toprule
\textbf{Rückgang} & Neue Anwendungen wählen überwiegend elliptische Kurven --- kürzere Schlüssel, schnellere Signaturen.\\
\textbf{Bestand} & In Zertifikaten und Altsystemen weiterhin allgegenwärtig.\\
\textbf{Quanten} & Shors Verfahren faktorisiert in polynomialer Zeit; längere Schlüssel helfen nicht.\\
\bottomrule
\end{tabular}}
\end{center}

\schritt{Zur ersten Zeile}
Bei gleicher Sicherheit ist eine elliptische Signatur etwa zwölfmal kürzer und ihre Erzeugung deutlich schneller.

\begin{center}
\begin{minipage}{0.88\textwidth}
\kursiv{RSA hat allerdings einen Vorteil, der oft übersehen wird: Die \kursiv{Prüfung} ist mit kleinem $e$ sehr schnell --- schneller als bei elliptischen Kurven. Wo eine Signatur einmal erzeugt und millionenfach geprüft wird, kann das den Ausschlag geben.}
\end{minipage}
\end{center}

\schritt{Zur dritten Zeile}
Shor veröffentlichte sein Verfahren 1994. Der Aufwand ist polynomial in der Bitlänge.

\begin{center}
\begin{minipage}{0.88\textwidth}
\kursiv{Nach heutigen Schätzungen bräuchte man für einen $2048$-Bit-Schlüssel einige Millionen fehlerkorrigierte Quantenbits --- gegenwärtig sind es einige hundert bis wenige tausend physikalische. \textbf{Der Abstand ist groß und schrumpft}, und niemand kann zuverlässig sagen, wie schnell. Kapitel~15 behandelt, was an die Stelle tritt.}
\end{minipage}
\end{center}
\end{forschung}

\begin{merksatz}
\boxtitel{Was Sie jetzt können}
\begin{center}
\renewcommand{\arraystretch}{1.5}
\begin{tabular}{@{}P{3.0cm}P{5.6cm}@{}}
\toprule
\textbf{Satz 13.1} & Euler: $a^{\varphi(n)}\equiv1$\\
\textbf{Satz 13.2} & $m^{ed} = m$ in vier Zeilen\\
\textbf{Die Sicherheit} & höchstens so hart wie Faktorisieren\\
\textbf{Der Bruchfall} & bestimmt, formbar, kleine Werte\\
\textbf{Signaturen} & den Hashwert signieren, nie die Nachricht\\
\bottomrule
\end{tabular}
\end{center}

\begin{center}
\begin{minipage}{0.88\textwidth}
\kursiv{Im nächsten Kapitel ersetzen wir die Restklassen durch eine Gruppe, in der dieselben Verfahren mit einem Zwölftel der Bits auskommen.}
\end{minipage}
\end{center}
\end{merksatz}
